% Generated draft: requires mathematical and source review.
\documentclass[11pt]{article}
\usepackage[margin=25mm]{geometry}
\usepackage{fontspec}
\setmainfont{Latin Modern Roman}
\usepackage{amsmath,amssymb,mathtools,unicode-math}
\setmathfont{Latin Modern Math}
\usepackage{enumitem,needspace,xurl,hyperref}
\setlength{\parindent}{0pt}
\setlength{\parskip}{5pt}
\setlength{\emergencystretch}{3em}
\newcommand{\N}{\mathbb{N}}
\newcommand{\Z}{\mathbb{Z}}
\newcommand{\Q}{\mathbb{Q}}
\newcommand{\R}{\mathbb{R}}
\newcommand{\C}{\mathbb{C}}
\newcommand{\F}{\mathbb{F}}
\newcommand{\A}{\mathbb{A}}
\newcommand{\PP}{\mathbb{P}}
\newcommand{\E}{\mathbb{E}}
\newcommand{\eps}{\varepsilon}
\newcommand{\dd}{\,\mathrm{d}}
\begin{document}
\section*{UnsolvedMath ID: 30000982}
\textbf{Source problem code:} catalogue entry
\subsection{Definitions and mathematical statement}
Let $G$ be a finite group. For a finite-rank free abstract group $F_d$ and an epimorphism $\pi:F_d\twoheadrightarrow G$, a discrete presentation with $s$ relations means
\[
\ker\pi=\langle\!\langle r_1,\ldots,r_s\rangle\!\rangle_{F_d}.
\]
Let $r(G)$ be the minimum such $s$ over finite presentations of $G$.

For the free profinite group $\widehat F_d$ and a continuous epimorphism $\widehat\pi:\widehat F_d\twoheadrightarrow G$, a profinite presentation with $s$ relations means that $\ker\widehat\pi$ is the \emph{closed} normal subgroup generated by $s$ elements of $\widehat F_d$. These relations may be profinite words. Let $\widehat r(G)$ be the minimum such $s$ over finite-rank profinite presentations.

The question is
\[
\forall G\text{ finite},\qquad r(G)=\widehat r(G)?
\]
Completion of a discrete presentation gives $\widehat r(G)\leq r(G)$, so the requested equality asks whether the profinite category can ever reduce the minimum number of relations. The minima are understood over presentations, as stated in the source, rather than a preselected generating tuple; a variant requiring a fixed generator number would need that additional convention. The relation count is normal generation of the kernel, not merely generation of a relation module as an abelian or group-ring module.
\subsection{Short English statement}
Can a finite group require fewer defining relations in the profinite category than in the discrete category? The question compares ordinary normal generation with closed normal generation, allowing profinite relators in the latter.
\subsection{Sources}
\begin{itemize}
\item Supplied report: \emph{Profinite and Asymptotic Group Theory} (2009), Oberwolfach Reports, DOI 10.4171/owr/2008/28, pp. 1537--1588; proposed year 2008. \url{https://doi.org/10.4171/owr/2008/28}.
\item Supplied Guralnick--Kantor--Kassabov--Lubotzky reference: \url{https://pages.uoregon.edu/kantor/PAPERS/GKKL4.pdf}.
\item Local UnsolvedMath registry, permanent ID 30000982, catalogue entry; created 25 August 2026, updated 16 September 2026. Its corrected\_verified note reports comparison with the official report. The original explicitly says the equality is unknown; the supplied 21 August 2026 triage labels it open and says later work still calls it a central problem. No PDF access or independent current-status check is asserted here.
\end{itemize}
\end{document}
